\documentclass[12pt,a4paper]{article}
\usepackage[margin=2.5cm]{geometry}
\usepackage{graphicx}
\usepackage{amsmath,amssymb}
\usepackage{booktabs}
\usepackage{hyperref}
\usepackage{float}

\title{Exclusive $J/\psi$ Electroproduction at CLAS12:\\
Monte Carlo Tuning with the Forward Tagger and NN Fast Simulation}
\author{V.~Kubarovsky}
\date{\today}

\begin{document}
\maketitle

\begin{abstract}
Tuning of the DIFFRAD Monte Carlo generator for exclusive $J/\psi$
electroproduction $ep \to e' p\, J/\psi$, $J/\psi \to e^+e^-$, with the
scattered electron in the CLAS12 Forward Tagger (quasi-photoproduction,
$Q^2 \lesssim 0.1$~GeV$^2$). The model uses a dipole $t$-dependence.
The combined CLAS12 RGA sample of 68 $J/\psi$ candidates
(52 at 10.6~GeV, 16 at 10.2~GeV) is acceptance-corrected with neural-network
fast simulation and fitted by MC reweighting.
\end{abstract}

\tableofcontents
\newpage

%% ==================================================================
\section{Cross-Section Model}

\begin{equation}
\frac{d\sigma_T}{dt} = \alpha_1 \left(1-\frac{W_{\rm th}^2}{W^2}\right)^{\alpha_2}
W^{\alpha_3} \left(1+\frac{Q^2}{m_\psi^2}\right)^{-\nu_T}
\cdot \frac{3\,(m_g^2 - t_{\rm min})^3}{(m_g^2-t)^4}
\end{equation}
\begin{equation}
\frac{d\sigma_L}{dt} = c_R\,\frac{Q^2}{m_\psi^2}\,\frac{d\sigma_T}{dt},
\qquad W_{\rm th} = M_N + m_\psi = 4.035~\text{GeV}.
\end{equation}
Here $t, t_{\rm min} < 0$ and the dipole factor is normalized:
$\int_{-\infty}^{t_{\rm min}} 3(m_g^2-t_{\rm min})^3/(m_g^2-t)^4\,dt = 1$,
so the $t$-integrated cross section is carried entirely by the $W$ and $Q^2$
factors. Near threshold $|t|_{\rm min}$ is large (2.7~GeV$^2$ at $W_{\rm th}$,
1.0~GeV$^2$ at $W=4.5$), which makes the correct $t_{\rm min}$ treatment
essential.

In Forward Tagger kinematics ($Q^2 \in [0.015, 0.1]$~GeV$^2$) the factors
$(1+Q^2/m_\psi^2)^{-\nu_T}$ and $c_R Q^2/m_\psi^2$ vary by less than 1\%:
the model has no usable $Q^2$ lever arm, and $\nu_T$, $c_R$, $\alpha_3$ are
kept fixed. The free parameters are $\alpha_2$ (threshold behaviour) and
$m_g^2$ (dipole $t$-slope).

%% ==================================================================
\section{Monte Carlo Generator}

\texttt{diffrad\_jpsi\_dipole.f90} (DIFFRAD family, one particle = one
generator). All model parameters ($\alpha_1,\alpha_2,\alpha_3,\nu_T,m_g^2,c_R$)
are read from the input file. The accept--reject proposal uses the same dipole
shape, so the sampling efficiency tracks $m_g^2$ automatically.

Two bugs were found and fixed during this work:
\begin{itemize}
\item \textbf{Dipole normalization sign}: the code evaluated
$(m_g^2 - |t_{\rm min}|)^3$ instead of $(m_g^2 + |t_{\rm min}|)^3$.
Since $|t|_{\rm min} > m_g^2$ below $W \approx 4.23$~GeV, the cross
section went negative and threshold events were silently rejected.
\item \textbf{32-bit overflow}: the attempt limit \texttt{maxtry = 100000*nev}
overflowed INTEGER*4 for $n_{\rm ev} \geq 2$M, terminating generation
instantly. Counters are now \texttt{integer*8} (also fixed in the $\phi$
generator).
\end{itemize}

Samples: 5M events at 10.6 GeV, 2M at 10.2 GeV.
Generation ranges: $Q^2 \in [0.01, 0.12]$, $W \in [4.05, 4.55]$,
$|t| < 6$~GeV$^2$, $\theta_{e'} \in [0^\circ, 7.5^\circ]$.

%% ==================================================================
\section{Data}

Combined RGA $J/\psi$ candidate sample (Mariana), 68 events
(52 at 10.6 GeV, 16 at 10.2 GeV), four detected particles:
$e'$ in the Forward Tagger ($\theta = 2.7$--$6.0^\circ$),
decay $e^+e^-$ and recoil proton in the Forward Detector.
$M(e^+e^-) \in [2.93, 3.17]$~GeV --- the sample is pre-selected, the yield
is obtained by counting (no mass fit). 64 events pass the analysis cuts
($Q^2 \in [0.015, 0.1]$,
$W \in [4.15, 4.5]$, $|t| < 3.0$).

%% ==================================================================
\section{Fast Monte Carlo and Acceptance}

NN models (acceptance + smearing) trained per particle and per run
configuration: $e^-$ (FT, no vertex measurement, $v_z$ fixed at $-3$~cm),
decay $e^-$, $e^+$, $p$. The 10.6 GeV MC is evaluated against all
8 run configurations (F18 in/out, S18 in/out at their beam
currents); each event's acceptance weight is the fraction of configurations
that accept it. The 10.2 GeV MC uses S19in. The two beam energies are then
combined so that the \emph{accepted} MC matches the data composition 52:16.

\begin{figure}[H]
\centering
\includegraphics[width=0.99\textwidth]{plots/step_gen_vs_acc.png}
\caption{Generated vs fastMC-accepted distributions for both beam energies.}
\end{figure}

%% ==================================================================
\section{Tuning Iterations}

Fit: acceptance-corrected counts in 3 $Q^2$ + 3 $W$ + 4 $|t|$ bins (10 points),
model predictions by event-by-event reweighting
$\sigma_{\rm new}/\sigma_{\rm gen}$, free parameters $\alpha_2$ and $m_g^2$,
\texttt{differential\_evolution} minimization.

\begin{table}[H]
\centering
\caption{Tuning iterations. Starting values $\alpha_2=1.0$, $m_g^2=1.6$
(Bhawani/lAger).}
\begin{tabular}{ccccccc}
\toprule
Iter & $\alpha_2^{\rm gen}$ & $m_g^{2,\rm gen}$ &
$\alpha_2^{\rm fit}$ & $m_g^{2,\rm fit}$ & $\chi^2$/ndf \\
\midrule
1 & 1.000 & 1.600 & 3.719 & 3.144 & 24.1/8 \\
2 & 3.719 & 3.144 & 4.048 & 3.152 & 17.9/8 \\
3 & 4.105 & 3.106 & 4.122 & 3.112 & 17.3/8 \\
\bottomrule
\end{tabular}
\end{table}

\begin{figure}[H]
\centering
\includegraphics[width=0.99\textwidth]{plots/jpsi_corrected.png}
\caption{Acceptance-corrected counts vs generation model (blue) and fitted
model (red). $Q^2$, $W$, $|t|$ are fitted; $|t|-|t|_{min}$ and $x_B$ are
cross-checks. The first $Q^2$ bin carries the dominant $\chi^2$ contribution;
the model $Q^2$ shape is flux-driven and not tunable in FT kinematics.}
\end{figure}

%% ==================================================================
\section{Cross-Section Shapes in the (W, Q$^2$) Analysis Bins}

Flux-folded $\gamma^* p \to J/\psi\,p$ shapes in the three correlated
$(W, Q^2)$ bins of the CLAS12 analysis (Table~16 of the analysis note).
Absolute normalization requires the integrated luminosity and is not
attempted here.

\begin{figure}[H]
\centering
\includegraphics[width=0.99\textwidth]{plots/jpsi_dsdt_mariana_bins.png}
\caption{$d\sigma/dt$ shape in the three $(W,Q^2)$ bins:
acceptance-corrected data (points) vs tuned model (red).}
\end{figure}

\begin{figure}[H]
\centering
\includegraphics[width=0.9\textwidth]{plots/jpsi_sigma_q2_w.png}
\caption{Cross-section shape vs $Q^2$ and $W$, integrated over $|t|<3$~GeV$^2$.}
\end{figure}

%% ==================================================================
\section{Acceptance Comparison with the CLAS12 Analysis}

\begin{table}[H]
\centering
\caption{Our NN-fastMC mixture acceptance vs the efficiencies of the CLAS12
analysis (Table~17): $\eta_1$ event-by-event, $\eta_{2a}$/$\eta_{2b}$
bin-averaged with two methods.}
\begin{tabular}{cccccccc}
\toprule
Bin & $Q^2$ [GeV$^2$] & $W$ [GeV] & $\varepsilon$ (this work) &
$\eta_1$ & $\eta_{2a}$ & $\eta_{2b}$ & $N_{\rm data}$ \\
\midrule
1 & [0.029, 0.071] & [4.202, 4.256] & 0.1016 & 0.1463 & 0.1184 & 0.1406 & 7 \\
2 & [0.018, 0.047] & [4.320, 4.387] & 0.1142 & 0.1101 & 0.1096 & 0.1386 & 11 \\
3 & [0.016, 0.043] & [4.426, 4.487] & 0.0869 & 0.1090 & 0.1127 & 0.1313 & 15 \\
\bottomrule
\end{tabular}
\end{table}

%% ==================================================================
\section{Summary}

The tuned $J/\psi$ generator parameters are $\alpha_2 = 4.12$, $m_g^2 = 3.11$~GeV$^2$ ($\alpha_3 = 0.32$, $\nu_T = 3.0$, $c_R = 0.4$ fixed), $\chi^2 = 17.3$ for 8~ndf. The data prefer a much steeper threshold rise and a flatter
$t$-dependence than the starting (photoproduction-inspired) parametrization.
With 68 events the two-parameter description is adequate; the remaining
tension is confined to the lowest $Q^2$ bin, where the model shape is fixed
by the virtual photon flux.

\end{document}
